\documentclass[11pt,a4paper]{article}

%% PACHETE MATEMATICE
\usepackage{amsmath,amssymb,amsfonts}
\usepackage{amsthm}
\usepackage{mathpazo}
\usepackage{eulervm}
\usepackage{enumerate}
\usepackage{color}
\usepackage{graphicx}
\usepackage[all]{xy}
\usepackage{xcolor}
\usepackage{framed}
\usepackage[unicode]{hyperref}
\setcounter{MaxMatrixCols}{20}
\usepackage{multicol}
% \usepackage[romanian]{babel}


%% ASPECT
\linespread{1.05}
\usepackage[T1]{fontenc}
\usepackage[utf8x]{inputenc}
\usepackage[margin=2cm]{geometry}
% \usepackage[color]{showkeys}
\usepackage{anyfontsize}
\usepackage{titlesec}
\titleformat*{\section}{\large\bfseries}

\usepackage{fancyhdr}
\pagestyle{fancy}
\rhead{\small \color{gray} Notițe de Adrian Manea}
\lhead{\small \color{gray} ETTI-AM2, 2017---2018, M. Joița \& A. Niță}


\usepackage[autostyle = false, style = german]{csquotes}
\usepackage[romanian]{babel}
\newcommand\qq{\enquote}

%% COMENZILE MELE
\renewcommand*{\proofname}{Dem.:}
\newcommand\bb{\mathbb}
\newcommand\dr{\mathrm}
\newcommand\kal{\mathcal}
\newcommand\fk{\mathfrak}
\newcommand\op{\oplus}
\newcommand\ot{\otimes}
\newcommand\ds{\displaystyle}
\newcommand\id{\indent}
\newcommand\nid{\noindent}
\newcommand\seq{\subseteq}
\newcommand\sneq{\subsetneq}
\newcommand\speq{\supseteq}
\newcommand\spneq{\supsetneq}
\newcommand\rar{\rightarrow}
\newcommand\Rar{\Rightarrow}
\newcommand\xrar{\xrightarrow}
\newcommand\lar{\leftarrow}
\newcommand\Lar{\Leftarrow}
\newcommand\xlar{\xleftarrow}
\newcommand\lrar{\leftrightarrow}
\newcommand\Lrar{\Leftrightarrow}
\newcommand\ideal{\trianglelefteq}
\newcommand\ai{\text{ a.\^{i}. }}
\newcommand\sk{(\textit{Schi\c{t}\u{a}}) }
\newcommand\rhup{\rightharpoonup}
\newcommand\rhdn{\rightharpoondown}
\newcommand\lhup{\leftharpoonup}
\newcommand\lhdn{\leftharpoondown}
\newcommand\vid{\emptyset}
\newcommand\wh{\widehat}
\newcommand\ol{\overline}
\newcommand\NN{\mathbb{N}}
\newcommand\RR{\mathbb{R}}
\newcommand\ZZ{\mathbb{Z}}
\newcommand\QQ{\mathbb{Q}}
\newcommand\CC{\mathbb{C}}

%% STILURILE MELE
\newtheoremstyle{dotless}{}{}{\itshape}{}{\bfseries}{:}{ }{}

\theoremstyle{dotless}
\newtheorem{theorem}{Teorem\u{a}}
\newtheorem{proposition}{Propozi\c{t}ie}
\newtheorem{lemma}{Lem\u{a}}[section]
\newtheorem{corollary}{Corolar}
\newtheorem{exercise}{Exerci\c{t}iu}

\newtheoremstyle{dotlessdr}{}{}{}{}{\bfseries}{:}{ }{}
\theoremstyle{dotlessdr}
\newtheorem{example}{Exemplu}
\newtheorem{definition}{Defini\c{t}ie}
\newtheorem{remark}{Observa\c{t}ie}




\begin{document}
% \definecolor{labelkey}{RGB}{222,0,0}

\begin{center}
  {\large\textbf{Seminar 8 \\
      Analiză complexă}}
\end{center}

\vspace{2cm}

\section*{Funcții complexe}

\indent\indent Următoarele noțiuni sînt elementare și introduc conceptele
de bază din studiul analizei complexe, începînd cu funcțiile de variabile complexe.

\begin{definition}\label{def:fc-comp}
  Fie $ A \seq \CC $. Se numește \emph{funcție complexă} orice funcție de forma
  $ f : A \to \CC $.
\end{definition}

Dată fiind scrierea algebrică a unui număr complex, putem separa părțile și pentru
funcții. Astfel, dacă $ w = f(z) $, iar $ z = x + iy \in A, w = u + iv \in \CC $,
putem scrie $ f = P + iQ $, cu $ P = \mathrm{Re}f$ și $ Q = \mathrm{Im}f $. Am pus
în evidență două funcții $ P, Q : A \seq \RR^2 \to \RR $, iar egalitatea $ w = f(z) $
devine echivalentă cu două egalități reale $ u = P(x, y), v = Q(x, y) $, funcția
în sine fiind echivalentă cu o transformare punctuală:
\[
  A \to \RR^2, \quad (x, y) \mapsto (P(x, y), Q(x, y)).
\]

De aceea, unele proprietăți le putem analiza pe componente:

\begin{definition}\label{def:fc-cont}
  Fie $ A \seq \CC $ o mulțime deschisă și $ f : A \to \CC $ o funcție complexă.

  Funcția este \emph{continuă} în punctul $ z_0 = x_0 + iy_0 $ dacă și numai
  dacă $ P $ și $ Q $ sînt simultan continue în punctul $ (x_0, y_0) $.
\end{definition}

Următoarea noțiune este legată de derivabilitate.

\begin{definition}\label{def:fc-der}
  Fie $ A \seq \CC $ o mulțime deschisă și $ f : A \to \CC $ o funcție complexă.

  Funcția $ f $ se numește \emph{olomorfă} în punctul $ z_0 \in A $ (echivalent,
  \emph{$\CC$-derivabilă} sau \emph{monogenă}) dacă există și este finită limita:
  \[
    \ell = \lim_{\substack{z \to z_0 \\ z \neq z_0}} \frac{f(z) - f(z_0)}{z - z_0}.
  \]

  În caz afirmativ, $ \ell = f'(z_0) $ se numește \emph{derivata complexă} a lui
  $ f $ în $ z_0 $.
\end{definition}

Similar cu cazul real, avem:

\begin{remark}\label{rk:ol-cont}
  Dacă funcția $ f : A \to \CC $ este olomorfă în $ z_0 \in A $, atunci ea este
  continuă în $ z_0 $.
\end{remark}

Noțiunile se pot extinde pentru întreg deschisul în mod evident: funcția se numește
\emph{continuă} (respectiv \emph{olomorfă}) pe $ A $ dacă are această proprietate în
orice punct din $ A $.

O condiție de derivabilitate care ține cont de descompunerea funcțiilor complexe
este următoarea:

\begin{theorem}\label{thm:ec-cauchy-riemann}
  Fie $ A \seq \CC $ o mulțime deschisă.

  Funcția $ f : A \to \CC, f = P + iQ $ este olomorfă în $ z_0 \in A $ dacă și numai
  dacă $ P, Q : A \to \RR $ sînt diferențiabile în $ z_0 = (x_0, y_0) $, iar derivatele
  lor parțiale în $ (x_0, y_0) $ verifică \emph{condițiile Cauchy-Riemann}, adică:
  \[
    \begin{cases}
      \frac{\partial P}{\partial x} &= \frac{\partial Q}{\partial y} \\
      \frac{\partial Q}{\partial x} &= -\frac{\partial P}{\partial y}
    \end{cases}
  \]
\end{theorem}

Următoarea noțiune ne ajută să găsim o condiție echivalentă cu ecuațiile
Cauchy-Riemann:

\begin{definition}\label{def:der-areolara}
  Dacă $ f : A \to \CC $ este o funcție complexă, definim:
  \[
    \begin{cases}
      \frac{\partial f}{\partial z} &= \frac{1}{2} \Big(\frac{\partial f}{\partial x} - %
      \frac{\partial f}{\partial y} \Big) \\
      \frac{\partial f}{\partial \overline{z}} &= \frac{1}{2}\Big( %
      \frac{\partial f}{\partial x} + i\frac{\partial f}{\partial y} \Big)
    \end{cases}
  \]

  Dintre acestea, $ \frac{\partial f}{\partial \overline{z}} $ se numește
  \emph{derivata areolară} a lui $ f $.
\end{definition}

Avem, atunci:

\begin{corollary}\label{cor:are-c-r}
  Relația $ \frac{\partial f}{\partial \overline{z}} = 0 $ este echivalentă cu
  condițiile Cauchy-Riemann.
\end{corollary}


Amintim din cazul funcțiilor reale:
\begin{definition}\label{def:fc-armonica}
  Fie $ u : A \to \RR $ o funcție de clasă $ \kal{C}^2 $ pe deschisul $ A $.
  Funcția $ u $ se numește \emph{armonică} dacă pentru orice punct $ a \in A $
  are loc:
  \[
    \frac{\partial^2 u}{\partial x^2}(a) + \frac{\partial^2 u}{\partial y^2}(a) = 0,
  \]
  adică $ \Delta u = 0 $ în orice punct $ a \in A $.
\end{definition}

Putem folosi această noțiune în următorul context, de exemplu:

\begin{corollary}\label{cor:olo-arm}
  Fie $ A \seq \CC $ o mulțime deschisă și $ f : A \to \CC, f = P + iQ $, cu
  $ P, Q \in \kal{C}^2(A) $.

  Dacă $ f $ este olomorfă pe $ A $, atunci $ P $ și $ Q $ sînt armonice pe $ A $.
\end{corollary}

Pentru rezultatul reciproc, avem nevoie de:

\begin{definition}\label{def:conex}
  O mulțime deschisă $ D \seq \CC $ se numește \emph{conexă} dacă pentru orice
  două puncte $ z_1, z_2 \in D $ există un drum $ \gamma: [a, b] \to D $
  care să unească cele două puncte, i.e.\ $ \gamma(a) = z_1, \gamma(b) = z_2 $.

  Domeniul se numește \emph{simplu conex} dacă frontiera lui este conexă.
\end{definition}

Cu aceasta, avem:

\begin{theorem}\label{thm:sconex-arm-olo}
  Fie $ D \seq \RR^2 $ un domeniu simplu conex.

  Dacă funcția $ P : D \to \RR $ este armonică pe $ D $, atunci există funcția
  $ Q : D \to \RR $ armonică pe $ D $ astfel încît funcția complexă
  $ f : D \to \CC, f = P + iQ $ să fie olomorfă.
\end{theorem}

%%%%%%%%%%%%%%%%%%%%%%%%%%%%%%%%%%%%%%%%%%%%%%%%%%%%%%%%%%%%%%%%%%%%%%

\section*{Funcții particulare}

\indent\indent În continuare, vom vedea cum funcții reale elementare, precum
funcția exponențială, radical, funcții trigonometrice etc.\ pot fi extinse pentru
a fi definite ca funcții complexe.

\begin{definition}\label{def:exp-c}
  Se numește \emph{exponențiala complexă} funcția $ \exp : \CC \to \CC $
  definită prin:
  \[
    \exp z = e^z = \sum_{n \geq 0} \frac{z^n}{n!}.
  \]
\end{definition}

Ca în cazul real, se pot demonstra ușor proprietățile:
\begin{enumerate}[(a)]
\item $ \exp(0) = 1 $;
\item $ \exp(z_1 + z_2) = \exp(z_1) \cdot \exp(z_2), \forall z_1, z_2 \in \CC $;
\item $ \exp(iy) = \cos y + i \sin y, \forall y \in \RR $ (Euler);
\item Funcția exponențială este olomorfă și periodică de perioadă $ T = 2\pi $.
\end{enumerate}

\vspace{1cm}

Pentru funcția logaritmică, dacă vrem să rezolvăm ecuația $ \exp(w) = z $, unde
$ w = u + iv \in \CC $, putem scrie în formă polară $ z = re^{i\theta} $ și
atunci, ținînd cont și de periodicitatea funcției exponențiale, găsim că:
\[
  \begin{cases}
    u &= \ln |z| \\
    v &= \mathrm{Arg}z + 2k\pi, k \in \ZZ
  \end{cases}
\]

Astfel, avem:

\begin{definition}\label{def:log-c}
  Se numește \emph{logaritmul numărului complex} $ z \in \CC^\ast $ mulțimea de
  numere complexe:
  \[
    \mathrm{Ln}z = \{ \ln|z| + i (\mathrm{Arg}z + 2k\pi) \mid k \in \ZZ \}.
  \]
\end{definition}

Această funcție este \emph{multiformă}, adică asociază unei valori $ z $ mai
multe valori numerice (o infinitate de ramuri, de fapt). Iar pentru $ k = 0 $,
obținem \emph{valoarea principală} a logaritmului:
\[
  \ln z = \ln |z| + i \mathrm{Arg}z.
\]

\vspace{1cm}

Putem acum defini și funcția putere și în particular, funcția radical:

\begin{definition}\label{def:putere-radical-c}
  Funcția putere de exponent complex:
  \[
    z^m = \exp(m \mathrm{Ln}z) = \big\{ \exp(m(\ln|z| + i(\mathrm{Arg}z + 2k\pi))) %
    \mid k \in \ZZ \big\}, m \in \CC.
  \]

  Funcția radical, de argument complex:
  \begin{align*}
    \sqrt[n]{z} &= z^{\frac{1}{n}} = \exp\big( \frac{1}{n}\mathrm{Ln}z \big) \\
                &= \big\{ \exp\big( \frac{1}{n}(\ln|z| + i(\mathrm{Arg}z + 2k\pi)) \big) \mid %
                  k \in \ZZ \big\} \\
                &= \Big\{ \exp\Big(\frac{1}{n} \ln|z|\Big) \cdot %
                  \exp\Big(i \frac{\theta + 2 k\pi}{n} \Big) \mid k \in \ZZ \} \\
                &= \Big\{ \sqrt[n]{r} \cdot \exp\Big( i \frac{\theta + 2 k\pi}{n}\Big) %
                  \mid k \in \NN \Big\}
  \end{align*}
\end{definition}

%%%%%%%%%%%%%%%%%%%%%%%%%%%%%%%%%%%%%%%%%%%%%%%%%%%%%%%%%%%%%%%%%%%%%%

Din funcția exponențială putem extrage și funcțiile trigonometrice complexe:

\begin{definition}\label{def:trig-c}
  Pentru orice $ z \in \CC $, definim:
  \begin{align*}
    \cos z &= \frac{e^{iz} + e^{-iz}}{2} \\
    \sin z &= \frac{e^{iz} - e^{-iz}}{2i} \\
    \tan z &= -i \frac{e^{iz} - e^{-iz}}{e^{iz} + e^{-iz}} \\
    \cot z &= i \frac{e^{iz} + e^{-iz}}{e^{iz} - e^{-iz}}
  \end{align*}
\end{definition}

\begin{remark}\label{rk:trig-unif}
  Funcțiile trigonometrice complexe sînt \emph{uniforme} (adică nu sînt multiforme),
  iar toate formulele din cazul real rămîn adevărate.
\end{remark}

Avem, de asemenea, și funcții trigonometrice hiperbolice:

\begin{definition}\label{def:trig-hip-c}
  \begin{align*}
    \sinh z &= \frac{e^z - e^{-z}}{2} \\
    \cosh z &= \frac{e^z + e^{-z}}{2} \\
    \tanh z &= \frac{\sinh z}{\cosh z}.
  \end{align*}
\end{definition}

Remarcăm că au loc legăturile:
\[
  \sinh z = -i\sin(iz), \quad \cosh z = \cos(iz).
\]

Rezolvarea unor ecuații trigonometrice ne conduce la introducerea
funcțiilor trigonometrice inverse:

\[
  z = \sin w \Rightarrow z = \frac{e^{iw} - e^{-iw}}{2i} \Leftrightarrow %
  e^{2iw} - 2ize^{iw} - 1 = 0,
\]
pe care o rezolvăm ca pe o ecuația de gradul al doilea și obținem:
\[
  w = \mathrm{Arcsin}z = -i\mathrm{Ln}(iz \pm \sqrt{1 - z^2}).
\]

Similar, obținem și:
\begin{align*}
  \mathrm{Arccos}z &= i\mathrm{Ln}(z \pm \sqrt{z^2 - 1}) \\
  \mathrm{Arctan}z &= -\frac{i}{2}\mathrm{Ln}\frac{i - z}{i + z}.
\end{align*}

%%%%%%%%%%%%%%%%%%%%%%%%%%%%%%%%%%%%%%%%%%%%%%%%%%%%%%%%%%%%%%%%%%%%%%

\section*{Exerciții}

1. Arătați că funcția $ f : \CC \to \CC, f(z) = |z| $ nu e olomorfă în
niciun punct din $ \CC $.

\emph{Indicație:} Pentru satisfacerea condițiilor Cauchy-Riemann, avem nevoie de
$ z = 0 $, dar în $ z = 0 $, partea reală a lui $ f $ nu are derivate parțiale.

\vspace{1cm}

2. Arătați că funcția:
\[
  f : \CC \to \CC, \quad f(z) = \sqrt{|z - \overline{z}^2}
\]
este continuă în $ z = 0 $, satisface condițiile Cauchy-Riemann în acest
punct, dar nu este olomorfă.

\emph{Indicație:} Partea reală a lui $ f $ nu este diferențiabilă în $ z = 0 $.
Într-adevăr, ar trebui să avem:
\[
  P(x, y) - P(0, 0) = 0 \cdot (x - 0) + 0 \cdot (y - 0) + P_1(z)|z - 0|,
\]
de unde $ \ds\lim_{z \to 0} P_1(z) = 0 $. Dar $ P_1(z) = \frac{P(z)}{|z|} $
și luînd două șiruri $ z = \frac{1}{n}, z' = \frac{1}{n} + i\frac{1}{n} $,
ambele tinzînd la $ 0 $, avem $ P_1(z_n) = 0 $, dar $ P_1(z'_n) = \sqrt{2} $,
deci $ P_1 $ nu are limită în origine.

\vspace{1cm}

3. Să se determine funcția olomorfă $ f = P + iQ $ pe $ \CC $, dacă
$ Q(x, y) = \varphi(x^2 - y^2), \varphi \in \kal{C}^2 $.

\emph{Soluție:} Fie $ \alpha = x^2 - y^2 $. Atunci:
\begin{align*}
  \frac{\partial Q}{\partial x} &= 2x\varphi'(\alpha) \\
  \frac{\partial Q}{\partial y} &= -2y\varphi'(\alpha) \\
  \Rightarrow \frac{\partial^2 Q}{\partial x^2} &= 2\varphi'(\alpha) + %
                                                  4x^2 \varphi^{\prime\prime}(\alpha) \\
  \frac{\partial^2 Q}{\partial y^2} &= -2\varphi'(\alpha) + 4y^2\varphi^{\prime\prime}(\alpha).
\end{align*}

Deoarece $ Q $ trebuie să fie armonică, avem $ \Delta Q = 0, \forall x, y $, de unde:
\[
  \varphi^{\prime\prime}(\alpha) = 0 \Rightarrow \varphi(\alpha) = c\alpha + c_1.
\]

Din condițiile Cauchy-Riemann pentru $ P $ și $ Q $, obținem acum:
\[
  \begin{cases}
    \frac{\partial P}{\partial x} = \frac{\partial Q}{\partial y} &= -2cy \\
    \frac{\partial P}{\partial y} = -\frac{\partial Q}{\partial x} &= -2cx
  \end{cases}
\]

Integrăm a doua ecuație și înlocuim în prima, pentru a obține:
\[
  P(x, y) = -2cxy + k.
\]

În fine:
\[
  f(z) = -2cxy + k + i(c(x^2 - y^2) + c_1) \Rightarrow f(z) = ciz^2 + d, c, d \in \RR.
\]

\vspace{1cm}

4. Fie $ P(x, y) = e^{2x} \cos 2y + y^2 - x^2 $. Să se determine funcția
olomorfă $ f = P + iQ $ pe $ \CC $ astfel încît $ f(0) = 1 $.

\emph{Soluție:} Verificăm că $ P $ este armonică. Verificăm condițiile
Cauchy-Riemann:
\begin{align*}
  \frac{\partial P}{\partial x} = \frac{\partial Q}{\partial y} &= 2e^{2x} \cos 2y - 2x \\
  -\frac{\partial P}{\partial y} = \frac{\partial Q}{\partial x} &= 2e^{2x} \sin 2y - 2y.
\end{align*}

Integrăm a doua ecuație în raport cu $ x $, înlocuim în prima și obținem:
\begin{align*}
  f(z) &= e^{2x} \cos 2y + y^2 - x^2 + i(e^{2x}\sin 2y - 2xy + k) \\
       &= e^{2x}(\cos 2y + i\sin 2y) - (x + iy)^2 + ki \\
  \Rightarrow f(z) &= e^{2z} - z^2 + ki.
\end{align*}

Folosind condiția din enunț, găsim $ k = 0 $.

\vspace{1cm}

5. Determinați soluțiile $ w \in \CC $ ale ecuației $ e^w = -2i $.

\vspace{1cm}

6. Rezolvați ecuația $ z^3 + 2 - 2i = 0 $.

\vspace{1cm}

7. Calculați:
\begin{enumerate}[(a)]
\item $ \sin(1 + i) $;
\item $ \sinh(1 - i) $;
\item $ \tan\big( \frac{\pi}{4} - i\ln 3 \big) $;
\item $ \tanh \big(\ln 2 + \frac{\pi i}{4}\big) $;
\item $ \mathrm{Arccos}(i \sqrt{3}) $.
\end{enumerate}

\vspace{1cm}

8. Rezolvați ecuația $ \sin z = 2 $.

\newpage

9. Fie funcțiile:
\begin{enumerate}[(a)]
\item $ f(z) = z \mathrm{Re}z $;
\item $ f(z) = z^2 + z \cdot \overline{z} + 3z - 2\overline{z} $.
\end{enumerate}

Determinați punctele în care $ f $ este derivabilă și să se calculeze $ f'(z) $
în aceste puncte.

\vspace{1cm}

10. Să se studieze olomorfia funcțiilor:
\begin{enumerate}[(a)]
\item $ f(z) = z $;
\item $ f(z) = \overline{z} $;
\item $ f(z) = \exp(z) $;
\item $ f(z) = \exp(\overline{z}) $;
\item $ f(z) = |z| $;
\item $ f(z) = 2z + z^2 $.
\end{enumerate}

\vspace{1cm}

11. Să se determine funcția olomorfă $ f(z) = u(x, y) + iv(x, y) $ dacă:
\begin{enumerate}[(a)]
\item $ u(x, y) = x^2 - y^2 - 2y $;
\item $ u(x, y) = x^4 - 6x^2 y^2 + y^4 $;
\item $ u(x, y) = (x \cos y - y\sin y)e^x $.
\end{enumerate}


\end{document}